% Píše: backend. Sekce kapitoly Testování a ověření (kapitolu otevírá testovani.tex).
\section{Testování backendu}
\label{sec:testovani-backend}

Backend se ověřoval ve třech vrstvách: automatickými testy na PC bez hardwaru, během v~režimu
simulace a~nakonec na skutečném Raspberry Pi se zapojenými součástkami.

\subsection{Automatické testy}

Testy jsou psané pro pytest~\cite{pytest} a~spouštějí se z~kořene repozitáře příkazem
\kod{python -m pytest}. Na PC s~Windows a~Pythonem 3.14 prošlo všech 746 testů za 2,1~s.
Rozdělení do souborů ukazuje tabulka~\ref{tab:testy}.

\begin{table}[H]
\centering
\caption{Automatické testy backendu}
\label{tab:testy}
\small
\begin{tabular}{|p{5.2cm}|p{1.3cm}|p{7.0cm}|}
\hline
\textbf{Soubory \kod{test_*.py}} & \textbf{Testů} & \textbf{Co ověřují} \\
\hline
\kod{nastaveni}, \kod{denni_doba}, \kod{termostat}, \kod{svetlo}, \kod{text_displeje} & 257 &
  čistou logiku: typy a~rozsahy nastavení, atomické uložení, úseky dne, tepelný model
  a~hysterezi, automatiku světla, řádky displeje \\
\hline
\kod{jadro}, \kod{simulace} & 164 & krok domku, stav, příkazy, poruchy a~kontrolku pohybu se
  simulovaným hardwarem \\
\hline
\kod{ble_format}, \kod{ble_server} & 128 & kódování ESS a~JSON, soulad dokumentu rozhraní
  s~kódem, server BLE s~falešnou knihovnou bluezero \\
\hline
\kod{ovladace_pi}, \kod{cidlo_tsl2561}, \kod{rgb_led}, \kod{systemovy_cas} & 105 & ovladače
  s~falešnými knihovnami, vzorec pro lux, gama korekci \\
\hline
\kod{spusteni}, \kod{casovac}, \kod{over_soucastky} & 92 & volby a~spuštění programu, obal
  \kod{opakuj}, zkoušky součástek \\
\hline
\textbf{celkem} & \textbf{746} & \\
\hline
\end{tabular}
\end{table}

Knihovny pro Raspberry Pi (smbus2, RPi.bme280, gpiozero, luma, bluezero, dbus a~gi) na PC
nejsou a~do vývojového prostředí se neinstalují. Soubor \kod{tests/conftest.py} proto obsahuje
jejich falešné náhrady, které fixtury \kod{falesny_hardware} a~\kod{falesny_bluetooth} vloží do
\kod{sys.modules} místo skutečných modulů. Náhrady mají stejné signatury jako skutečné knihovny
ve verzích, jejichž zdrojový kód se při psaní ovladačů četl (smbus2 0.6.1, RPi.bme280 0.2.4,
gpiozero 2.0.1, luma.core 2.6.0, luma.oled 3.16.0, bluezero 0.9.1 a~dbus-python 1.3.2), takže
ovladač, který by je volal špatně, selže i~v~testu. Zaznamenávají, co ovladač volal, a~umějí na
zvoleném místě selhat (otevření sběrnice, čtení ID, měření, zápis na GPIO, inicializace displeje),
aby šlo ověřit hlášení poruch. Falešná bluezero napodobuje registraci, čtení s~posunem, zápis,
odběr i~zánik registrace v~BlueZ.

Testy jsou několika druhů:
\begin{itemize}
  \item jednotkové testy čisté logiky (nastavení, denní doba, termostat, světlo, text displeje,
        formát BLE), včetně hraničních hodnot a~pořadí kroků,
  \item testy jádra se simulovaným hardwarem z~\kod{hardware/simulace.py}, jehož atribut
        \kod{selhava} napodobí poruchu součástky,
  \item testy ovladačů a~serveru BLE s~falešnými knihovnami,
  \item testy spuštění, které rozběhnou program celý, v~simulaci i~v~režimu Raspberry Pi
        s~falešným hardwarem a~Bluetooth, s~falešnými hodinami místo čekání,
  \item testy dokumentu \kod{docs/ble-rozhrani.md}: UUID, klíče Stavu, výchozí hodnoty
        a~rozsahy Nastavení, názvy poruch a~všechny příklady JSON musí odpovídat kódu.
\end{itemize}

\subsection{Kontrola změn}

Každou ucelenou změnu kódu po dokončení zkontrolovala nezávislá recenze: recenzent, který změnu
nepsal, ji porovnal se zadáním a~posoudil kvalitu kódu a~testů. Nálezy se třídily na
kritické, důležité a~drobné; kritické a~důležité se opravily a~oprava prošla novou recenzí.
Nakonec prošla recenzí i~celá sada změn najednou.

Sílu testů ověřovaly mutační testy. Do kopie kódu se vždy vložila jedna úmyslná chyba (mutant),
například obrácená nerovnost, vynechaná podmínka nebo jiná konstanta, a~celá sada testů se
spustila znovu. Mutant, kterého žádný test nezachytil, ukazuje na chybějící test. Při vzniku
modulu nastavení testy zachytily všech 16 mutantů a~u~denní doby s~termostatem všech 44.
Recenze kontrolky pohybu takto našla chybějící test: mutant, který při trvajícím pohybu
přestal prodlužovat dobu svícení, testy nezachytily, a~proto přibyl test, který ho zachytí.

\subsection{Simulace na PC}

Režim \kod{--simulace} (oddíl~\ref{sec:simulace}) rozběhne celý domek na PC se simulovanými
čidly a~každou sekundu vypíše stav. Ověřuje se v~něm chování v~čase: že vnitřní teplota
stoupá k~cíli a~termostat se po jeho dosažení vypíná, že úsek dne odpovídá času a~že stav má
stejné klíče jako Stav v~BLE.

\subsection{Ověření na hardwaru}

Na Raspberry Pi se nejdřív každá součástka vyzkoušela zvlášť zkouškami \kod{over_soucastky}
(oddíl~\ref{sec:simulace}) a~potom celá služba. Zaznamenané výsledky z~21.~až 23.~10.~2026
shrnuje tabulka~\ref{tab:hw-vysledky}.

\begin{table}[H]
\centering
\caption{Zkoušky na hardwaru a~jejich výsledky}
\label{tab:hw-vysledky}
\small
\begin{tabular}{|p{2.6cm}|p{4.8cm}|p{6.1cm}|}
\hline
\textbf{Zkouška} & \textbf{Co dělá} & \textbf{Výsledek} \\
\hline
\kod{i2cdetect -y 1} & vypíše adresy na sběrnici I\textsuperscript{2}C & \kod{39} (TSL2561),
  \kod{3c} (SSD1306) a~\kod{76} (BMP280) \\
\hline
\kod{bme280} & ID čipu z~registru \kod{0xD0} a~tři měření po sekundě & ID \kod{0x58} (BMP280),
  23,3~°C, 985,1~hPa, vlhkost \kod{--} \\
\hline
\kod{tsl2561} & ID čipu z~registru \kod{0x8A} a~deset měření kanálů CH0, CH1 a~lux po
  sekundě & ID \kod{0x50}; CH0~=~62 a~CH1~=~9, tedy 26,0~lx; večer 4,5~lx \\
\hline
\kod{pir} & 30~s vypisuje změny čidla pohybu & výsledek nezaznamenaný; výstup čidla před
  připojením změřen multimetrem: 3,3~V \\
\hline
\kod{led} & třikrát blikne červená, zelená a~oranžová LED & výsledek nezaznamenaný \\
\hline
\kod{displej} & na 5~s ukáže testovací text se všemi šesti řádky & rozložení řádků zkontrolované
  náhledem s~písmem knihovny (obrázek~\ref{fig:displej}) \\
\hline
\kod{rgb} & rozsvítí červenou, zelenou, modrou a~bílou, každou na 1~s & pořadí barev správné;
  večerní barva po gama korekci oranžová \\
\hline
služba \kod{domek} & běh celého programu & \kod{active}, bez restartů, pole \kod{chyby} prázdné,
  čas z~NTP, periferie \kod{Termostat} zaregistrovaná v~BlueZ \\
\hline
uložené nastavení & načtení souboru z~dřívější verze & starý klíč převeden, zůstaly hodnoty
  nastavené z~tabletu, například barva $(46, 197, 232)$ a~jas 44~\% \\
\hline
\end{tabular}
\end{table}

Uložené nastavení zároveň dokládá, že zápisy z~tabletu přes BLE na hardwaru fungují: barvu,
jas, práh šera a~začátek večera v~souboru nastavil tablet. Pro ruční zkoušku rozhraní je
v~README backendu doporučený postup s~obecnou aplikací nRF Connect for
Mobile~\cite{nrfconnect}: připojit se, vyžádat MTU 517, zapnout notifikace Stavu, zapsat platnou
i~neplatnou hodnotu do Nastavení, poslat příkaz světla, odpojit se a~znovu připojit a~přečíst
Stav bez zvýšení MTU. Výsledky tohoto postupu po jednotlivých krocích zaznamenané nejsou.

\subsection{Problémy nalezené na skutečném hardwaru}
\label{sec:problemy-hw}

Testy na PC pracují s~falešnými knihovnami, které odpovídají zdrojovému kódu knihoven, ale ne
nutně skutečným součástkám. Na hardwaru se proto objevilo pět problémů, které testy zachytit
nemohly:
\begin{description}
  \item[BMP280 posílal vlhkost 0~\%.] Knihovna RPi.bme280 u~čidla BMP280, které vlhkost neměří,
        vracela vlhkost nulovou a~tablet by ukazoval 0~\%. Ovladač teď podle ID čipu \kod{0x58}
        pozná BMP280 a~vlhkost hlásí jako neznámou (\kod{null}, v~ESS \kod{0xFFFF}) bez
        poruchy. Displej místo vlhkosti ukazuje tlak a~zkouška \kod{bme280} na BMP280
        upozorní.
  \item[Čidlo osvětlení bylo TSL2561, ne BH1750.] Původní program počítal s~BH1750 na adrese
        \kod{0x23}, ale na sběrnici se ozvalo čidlo na adrese \kod{0x39} s~ID \kod{0x50}, tedy
        TSL2561. Program dostal vlastní ovladač podle katalogového listu~\cite{tsl2561}
        (oddíl~\ref{sec:ovladace}) a~ovladač BH1750 se odstranil.
  \item[Večerní barva byla bílá.] Teplá žlutá $(255, 147, 41)$ s~jasem 60~\% svítila na RGB LED
        skoro bíle, protože střída PWM rostla se složkou barvy lineárně. Ovladač teď používá gama
        korekci 2,2 (oddíl~\ref{sec:svetlo}); nejdřív se upravily testy, které na původním kódu
        selhaly, a~po opravě prošly. Na hardwaru pak světlo svítilo oranžově.
  \item[Řádky displeje na rozhraní barev.] Řádky byly původně po 10 bodech, takže druhý řádek
        zasahoval do žlutého pruhu a~ležel na rozhraní barev. Rozložení se změnilo tak, že první
        řádek je sám ve žlutém pruhu, a~potom se ještě posunulo na polohy 0, 14, 24, 33, 42
        a~52 bodů, aby se ocásky písmen nedotýkaly dalšího řádku.
  \item[Prázdný config.txt po nečistém vypnutí.] Raspberry Pi přišlo o~napájení krátce po
        instalaci, při které \kod{raspi-config} zapsal zapnutí I\textsuperscript{2}C do
        \kod{/boot/firmware/config.txt}. Soubor na oddílu FAT zůstal prázdný (0~bajtů),
        I\textsuperscript{2}C se po startu nezapnulo a~domek hlásil poruchy \kod{meteo},
        \kod{osvetleni} i~\kod{displej}; v~logu byla chyba
        \kod{No such file or directory: /dev/i2c-1}. Oprava byla
        \kod{sudo raspi-config nonint do_i2c 0}, \kod{sync} a~restart. Instalační skript od té
        doby přepisuje \kod{config.txt} jen tehdy, když I\textsuperscript{2}C zapnuté není,
        a~hned potom zavolá \kod{sync}. Raspberry Pi se má vypínat příkazem
        \kod{sudo poweroff}, ne vytažením napájení.
\end{description}
